\begin{figure*}[t]
\centering
\includegraphics[width=.97\textwidth]{figures/observed_continuation_mechanism.pdf}
\caption{Observed-space dependence of terminal continuation across three validation-map pairs. Within each pair, the complete aircraft state, goal, primary rollout and 15-control successor bank are identical. P69 is evaluated in four obstacle maps (a,b,d,e), and P246 in two distinct forest layouts (g,h). All maps accumulate five lidar scans up to $t=0.8$ s. The background shows two-metre voxels in the slice $z\in[20,22)$ m; curves are horizontal projections of three-dimensional rollouts. The valid forest continuation in (g) ends at $z_{\mathrm{end}}=11.6$ m. Crosses mark rejected rollout endpoints. Panels (c,f,i) report all 15 binary control outcomes, the candidate witness label, successor integration steps until the first witness or bank exhaustion, and the number of changed observation features. All feasibility outcomes are computed by the three-dimensional validator.}
\label{fig:observed_mechanism}
\end{figure*}

The paired-map experiments in Fig.~\ref{fig:observed_mechanism} isolate the dependence of terminal continuation on observed space. Across six distinct validation scenes, P69 changes from no valid successor to one valid successor in each of the first two pairs, while P246 changes from one valid successor to none in the forest pair. Every primary remains feasible in both members of its pair. The complete state, goal, primary trajectory, successor controls and all 28 state--task features remain unchanged within each pair, whereas 18, 18 and 12 of the 30 observation features change. A state-only predictor therefore assigns the same input representation to opposite continuation labels; observation conditioning supplies the geometric information needed to distinguish them. These mechanism examples are drawn from the existing validation set of 767 inputs, including 517 Full-positive inputs.
